\section*{ID9964: Laboratory-Scale Experiments: δφ ∼ 10\textasciicircum{}-15}
\paragraph{Metadata.}
PiID: 8922569596881 | RTEID: RTE:P38692.4248:T0.9214 | UUID: 11562bb2-85ef-54a2-b3b7-f3cb301a28ef

\paragraph{Canonical Statement.}
The phase shift in entangled photon pairs traversing gravitational fields provides the most promising near-term test. Current interferometric precision approaches the required sensitivity of \$\textbackslash{}delta\textbackslash{}phi \textbackslash{}sim 10\textasciicircum{}\{-15\}\$ radians for kilometer-scale baselines with entangled optical photons.

\paragraph{Canonical Equation.}
\[
\delta\phi \sim 10^{-15}
\]

\paragraph{Dictionary.}
The phase shift in entangled photon pairs traversing gravitational fields provides the most promising near-term test.

\paragraph{Encyclopedia.}
Laboratory-Scale Experiments: δφ ∼ 10\textasciicircum{}-15 is a scaling / order-of-magnitude relation, written δφ ∼ 10\textasciicircum{}-15. It expresses a scaling or proportionality, capturing how the quantity grows with its parameters. It is built from δ, φ. Within the Trawin Zero Point registry it serves as one element of the framework's mathematical narrative.
